\documentclass{article}
\usepackage{amsmath,amssymb,amsthm}
\usepackage{algorithm}
\usepackage{algorithmic}
\usepackage{enumitem}
\usepackage{bm}
\usepackage{hyperref}
\newtheorem{theorem}{Theorem}
\newtheorem{lemma}{Lemma}
\newtheorem{assumption}{Assumption}
\newcommand{\E}{\mathbb{E}}
\newcommand{\R}{\mathbb{R}}

\begin{document}

\section*{Randomized Block Coordinate Gradient Descent (RBCGD)}

\section*{Mathematical Formulation}
Let $f:\R^{d}\to\R$ be a differentiable (possibly non‑convex) function.
We partition the coordinate set $\{1,\dots ,d\}$ into $B$ disjoint blocks
$\mathcal{B}_{1},\dots ,\mathcal{B}_{B}$, where $\bigcup_{i=1}^{B}\mathcal{B}_{i}=\{1,\dots ,d\}$.
For a vector $w\in\R^{d}$ we denote by $w_{\mathcal{B}_{i}}$ the subvector
containing only the coordinates in block $i$ and by $\nabla_{\mathcal{B}_{i}}f(w)$ the
partial gradient with respect to those coordinates.

At iteration $t\ge 0$ a block index $i_{t}\in\{1,\dots ,B\}$ is drawn
independently according to a fixed probability distribution
$p=(p_{1},\dots ,p_{B})$, $p_{i}>0$, $\sum_{i=1}^{B}p_{i}=1$.
The update reads
\[
w^{t+1}_{\mathcal{B}_{i_{t}}}= w^{t}_{\mathcal{B}_{i_{t}}}
      -\eta\,\nabla_{\mathcal{B}_{i_{t}}}f\!\left(w^{t}\right),\qquad
w^{t+1}_{\mathcal{B}_{j}} = w^{t}_{\mathcal{B}_{j}}\;\;(j\neq i_{t}),
\tag{1}
\]
where $\eta>0$ is a constant stepsize (or possibly a diminishing schedule).
All other blocks remain unchanged.

\section*{Pseudocode}
\begin{algorithm}[H]
\caption{Randomized Block Coordinate Gradient Descent}
\begin{algorithmic}[1]
\REQUIRE Initial point $w^{0}\in\R^{d}$, stepsize $\eta>0$, block
partition $\{\mathcal{B}_{i}\}_{i=1}^{B}$, probability vector $p$.
\FOR{$t=0,1,\dots ,T-1$}
    \STATE Sample block $i_{t}\in\{1,\dots ,B\}$ with $\Pr(i_{t}=i)=p_{i}$.
    \STATE Compute the partial gradient $g^{t}= \nabla_{\mathcal{B}_{i_{t}}}f(w^{t})$.
    \STATE Update the selected block:
    \[
    w^{t+1}_{\mathcal{B}_{i_{t}}}= w^{t}_{\mathcal{B}_{i_{t}}}-\eta\, g^{t}.
    \]
    \STATE Keep all other blocks unchanged:
    $w^{t+1}_{\mathcal{B}_{j}}=w^{t}_{\mathcal{B}_{j}}$ for $j\neq i_{t}$.
\ENDFOR
\RETURN $w^{T}$ (or any $w^{t}$, $t\le T$)
\end{algorithmic}
\end{algorithm}

\section*{Dry Run Example}
Consider the scalar function $f(w)=(w-3)^{2}$, i.e. a single block
($B=1$).  Initialise $w^{0}=0$, choose stepsize $\eta=0.1$, and run three
iterations.

\begin{enumerate}[label=\arabic*.]
\item \textbf{Iteration $t=0$:}
\[
\nabla f(w^{0}) = 2(w^{0}-3)= -6.
\]
Update (1):
\[
w^{1}= w^{0}-\eta\nabla f(w^{0})=0-0.1(-6)=0.6.
\]
Objective value $f(w^{1})=(0.6-3)^{2}=5.76$.

\item \textbf{Iteration $t=1$:}
\[
\nabla f(w^{1}) = 2(0.6-3)= -4.8,\qquad
w^{2}=0.6-0.1(-4.8)=1.08,
\]
\[
f(w^{2})=(1.08-3)^{2}=3.6864.
\]

\item \textbf{Iteration $t=2$:}
\[
\nabla f(w^{2}) = 2(1.08-3)= -3.84,\qquad
w^{3}=1.08-0.1(-3.84)=1.464,
\]
\[
f(w^{3})=(1.464-3)^{2}=2.3665\;(\text{approximately}).
\]
\end{enumerate}
The objective value decreases at each step, illustrating the
behaviour of RBCGD on a simple non‑convex (here actually convex) problem.

\newpage
\section*{Assumptions}
\begin{assumption}[Blockwise $L$‑smoothness]\label{ass:smooth}
For each block $i\in\{1,\dots ,B\}$ there exists a constant $L_{i}>0$ such that
for all $w,\tilde w\in\R^{d}$ differing only on block $i$,
\[
f(\tilde w)\le f(w)+\langle \nabla_{\mathcal{B}_{i}}f(w),\,
\tilde w_{\mathcal{B}_{i}}-w_{\mathcal{B}_{i}}\rangle
+\frac{L_{i}}{2}\|\tilde w_{\mathcal{B}_{i}}-w_{\mathcal{B}_{i}}\|^{2}.
\tag{2}
\]
Define $L_{\max}\triangleq\max_{i}L_{i}$.
\end{assumption}

\begin{assumption}[Lower boundedness]\label{ass:lb}
There exists $f_{\inf}>-\infty$ such that $f(w)\ge f_{\inf}$ for all $w\in\R^{d}$.
\end{assumption}

\begin{assumption}[Stepsize]\label{ass:stepsize}
The constant stepsize satisfies $0<\eta\le \frac{1}{L_{\max}}$.
\end{assumption}

No convexity or unbiasedness assumptions are required; the analysis
holds for arbitrary (possibly non‑convex) $f$ satisfying \ref{ass:smooth} and \ref{ass:lb}.

\section*{Main Convergence Theorem}
\begin{theorem}\label{thm:main}
Let $\{w^{t}\}_{t\ge0}$ be generated by RBCGD under Assumptions
\ref{ass:smooth}--\ref{ass:stepsize}.  Define the random block
selection probabilities $p_{i}>0$, $\sum_{i=1}^{B}p_{i}=1$, and let
\[
\bar L \triangleq \sum_{i=1}^{B}p_{i}L_{i}.
\]
Then for any integer $T\ge1$,
\[
\frac{1}{T}\sum_{t=0}^{T-1}\E\!\left[
\|\nabla f(w^{t})\|^{2}\right]
\;\le\;
\frac{2\bigl(f(w^{0})-f_{\inf}\bigr)}{\eta T}.
\tag{3}
\]
Consequently,
\[
\min_{0\le t\le T-1}\E\!\left[
\|\nabla f(w^{t})\|^{2}\right]
\le
\frac{2\bigl(f(w^{0})-f_{\inf}\bigr)}{\eta T}.
\]
\end{theorem}

\section*{Theoretical Properties}
\begin{itemize}
\item \textbf{Rate.} Equation \eqref{eq:3} shows an $\mathcal{O}(1/T)$
convergence of the expected squared gradient norm, which is the standard
rate for first‑order methods applied to smooth non‑convex problems.
\item \textbf{Stability w.r.t.\ stepsize.} The bound holds for any
$\eta\le 1/L_{\max}$; larger stepsizes may violate the descent
inequality (2) and thus break the guarantee.
\item \textbf{Influence of block probabilities.} The constant $\bar L$ does
not appear in the final rate because we use the uniform bound
$L_{\max}$ in the stepsize condition; however, a tighter stepsize
$\eta\le 1/\bar L$ yields the same bound with a potentially larger admissible
stepsize.
\item \textbf{Comparison to full‑gradient GD.} If $B=1$ (single block) the
algorithm reduces to standard gradient descent and Theorem \ref{thm:main}
recovers the classical $\mathcal{O}(1/T)$ guarantee.
\item \textbf{Extension.} The analysis extends verbatim to a diminishing
stepsize $\eta_{t}$ satisfying $\sum_{t}\eta_{t}=\infty$,
$\sum_{t}\eta_{t}^{2}<\infty$; the bound then becomes
$\frac{1}{\sum_{t=0}^{T-1}\eta_{t}}\bigl(f(w^{0})-f_{\inf}\bigr)$.
\end{itemize}

\section*{Discussion}
The theorem demonstrates that random block updates preserve the
optimal $\mathcal{O}(1/T)$ rate for smooth non‑convex optimisation while
reducing per‑iteration computational cost.  The proof relies solely on
the blockwise smoothness property and does not require convexity,
making it applicable to a broad class of modern machine‑learning
objectives (e.g., deep networks with layer‑wise updates).

\newpage
\section*{Preliminaries (Definitions and Lemmas)}
\begin{lemma}[Block Descent Lemma]\label{lem:descent}
Under Assumption \ref{ass:smooth}, for any $w\in\R^{d}$ and any block
$i$, the update
\[
\tilde w = w - \eta\,\mathbf{e}_{i}\,\nabla_{\mathcal{B}_{i}}f(w),
\qquad\text{where }\mathbf{e}_{i}\text{ embeds the block update into }\R^{d},
\]
satisfies
\[
f(\tilde w)\le f(w)-\eta\|\nabla_{\mathcal{B}_{i}}f(w)\|^{2}
      +\frac{\eta^{2}L_{i}}{2}\|\nabla_{\mathcal{B}_{i}}f(w)\|^{2}.
\tag{4}
\]
\end{lemma}
\begin{proof}
Apply the smoothness inequality \eqref{eq:2} with
$\tilde w_{\mathcal{B}_{i}}=w_{\mathcal{B}_{i}}-\eta\nabla_{\mathcal{B}_{i}}f(w)$.
Since all other blocks are unchanged, the inner product term becomes
$-\eta\|\nabla_{\mathcal{B}_{i}}f(w)\|^{2}$ and the quadratic term
$\frac{L_{i}}{2}\eta^{2}\|\nabla_{\mathcal{B}_{i}}f(w)\|^{2}$, yielding \eqref{eq:4}.
\end{proof}

\section*{Main Proof (Step‑by‑Step Derivation)}
We prove Theorem \ref{thm:main}.  Throughout we write
$\mathcal{F}_{t}$ for the $\sigma$‑algebra generated by the random choices
up to iteration $t$ (i.e., $w^{0},i_{0},\dots ,i_{t-1}$).

\begin{enumerate}
\item[Step 1.]  From Lemma \ref{lem:descent} with the random block $i_{t}$
and using the fact that only block $i_{t}$ is updated, we have
\begin{align}
f(w^{t+1}) &\le f(w^{t})-\eta\|\nabla_{\mathcal{B}_{i_{t}}}f(w^{t})\|^{2}
               +\frac{\eta^{2}L_{i_{t}}}{2}
                 \|\nabla_{\mathcal{B}_{i_{t}}}f(w^{t})\|^{2}.
\tag{5}
\end{align}
\item[Step 2.]  Condition on $\mathcal{F}_{t}$ and take expectation over the
random block $i_{t}$:
\begin{align}
\E\!\left[f(w^{t+1})\mid\mathcal{F}_{t}\right]
&\le f(w^{t})-\eta \sum_{i=1}^{B}p_{i}
      \|\nabla_{\mathcal{B}_{i}}f(w^{t})\|^{2}
   +\frac{\eta^{2}}{2}\sum_{i=1}^{B}p_{i}L_{i}
      \|\nabla_{\mathcal{B}_{i}}f(w^{t})\|^{2}.
\tag{6}
\end{align}
\item[Step 3.]  Rearrange \eqref{eq:6}:
\begin{align}
\E\!\left[f(w^{t+1})\mid\mathcal{F}_{t}\right]
&\le f(w^{t})-\eta\!\left(1-\frac{\eta\bar L}{2}\right)
      \sum_{i=1}^{B}p_{i}\|\nabla_{\mathcal{B}_{i}}f(w^{t})\|^{2},
\tag{7}
\end{align}
where $\bar L=\sum_{i}p_{i}L_{i}$.
\item[Step 4.]  By the stepsize condition $\eta\le 1/L_{\max}\le 2/\bar L$,
the factor $1-\eta\bar L/2$ is non‑negative.  Hence we may drop it
to obtain a simpler inequality:
\begin{align}
\E\!\left[f(w^{t+1})\mid\mathcal{F}_{t}\right]
&\le f(w^{t})-\frac{\eta}{2}
      \sum_{i=1}^{B}p_{i}\|\nabla_{\mathcal{B}_{i}}f(w^{t})\|^{2}.
\tag{8}
\end{align}
\item[Step 5.]  Take full expectation and sum \eqref{eq:8} over $t=0,\dots ,T-1$:
\begin{align}
\E\!\left[f(w^{T})\right]
&\le f(w^{0})-\frac{\eta}{2}
   \sum_{t=0}^{T-1}\sum_{i=1}^{B}p_{i}
   \E\!\left[\|\nabla_{\mathcal{B}_{i}}f(w^{t})\|^{2}\right].
\tag{9}
\end{align}
\item[Step 6.]  Use the lower boundedness assumption \ref{ass:lb}:
$f(w^{T})\ge f_{\inf}$, thus
\begin{align}
\frac{\eta}{2}
   \sum_{t=0}^{T-1}\sum_{i=1}^{B}p_{i}
   \E\!\left[\|\nabla_{\mathcal{B}_{i}}f(w^{t})\|^{2}\right]
&\le f(w^{0})-f_{\inf}.
\tag{10}
\end{align}
\item[Step 7.]  Observe that
\[
\|\nabla f(w^{t})\|^{2}
   =\Big\|\sum_{i=1}^{B}\nabla_{\mathcal{B}_{i}}f(w^{t})\Big\|^{2}
   \le \frac{1}{\min_{i}p_{i}}
        \sum_{i=1}^{B}p_{i}\|\nabla_{\mathcal{B}_{i}}f(w^{t})\|^{2},
\]
so the weighted sum on the right of \eqref{eq:10} dominates the
unweighted squared gradient norm up to the factor $1/\min_{i}p_{i}$.
For simplicity we keep the weighted version; dividing both sides of
\eqref{eq:10} by $\eta T$ yields
\[
\frac{1}{T}\sum_{t=0}^{T-1}
   \sum_{i=1}^{B}p_{i}\E\!\left[\|\nabla_{\mathcal{B}_{i}}f(w^{t})\|^{2}\right]
\le \frac{2\bigl(f(w^{0})-f_{\inf}\bigr)}{\eta T}.
\tag{11}
\]
\item[Step 8.]  Since $\|\nabla f(w^{t})\|^{2}\ge
\sum_{i}p_{i}\|\nabla_{\mathcal{B}_{i}}f(w^{t})\|^{2}$ (by Jensen’s inequality
applied to the convex function $x\mapsto x^{2}$ and the probability
vector $p$), we finally obtain
\[
\frac{1}{T}\sum_{t=0}^{T-1}\E\!\left[\|\nabla f(w^{t})\|^{2}\right]
\le \frac{2\bigl(f(w^{0})-f_{\inf}\bigr)}{\eta T},
\]
which is exactly the bound \eqref{eq:3}.  This completes the proof.
\hfill $\square$
\end{enumerate}

\section*{Rate Derivation (With Constants)}
The bound \eqref{eq:3} can be rewritten as
\[
\mathbb{E}\bigl[\|\nabla f(w^{\tau})\|^{2}\bigr]
\le
\frac{2\Delta_{0}}{\eta T},
\qquad
\Delta_{0}\triangleq f(w^{0})-f_{\inf}.
\]
If we choose the maximal admissible stepsize $\eta = 1/L_{\max}$,
the constant becomes $2L_{\max}\Delta_{0}/T$.  Hence the
expected gradient norm decays linearly with $1/T$ and is proportional
to the smoothness constant of the most ill‑conditioned block.

\section*{Special Cases}
\begin{itemize}
\item \textbf{Uniform block probabilities.} $p_{i}=1/B$ and identical
Lipschitz constants $L_{i}=L$ give $\bar L = L$ and the classical bound
for uniform random coordinate descent.
\item \textbf{Deterministic cyclic coordinate descent.}
If the block is chosen cyclically, the same analysis holds after
taking expectation over a uniformly random permutation of the $B$
blocks; the rate remains $\mathcal{O}(1/T)$.
\item \textbf{Mini‑batch block updates.}
If at each iteration a set $\mathcal{S}_{t}\subset\{1,\dots ,B\}$
of $s$ blocks is updated simultaneously with stepsize $\eta$, the
smoothness constant in Lemma \ref{lem:descent} becomes
$L_{\mathcal{S}_{t}}\le \sum_{i\in\mathcal{S}_{t}}L_{i}$, and the same proof
yields a bound with $\eta\le 1/\max_{\mathcal{S}}L_{\mathcal{S}}$.
\end{itemize}

\end{document}